\section{Application Architecture}\label{sec:studio-architecture}

Motivo Studio is a browser-hosted IDE for the Motivo language.
The user edits \texttt{.motivo} source, compiles it through a backend API, inspects compiler diagnostics in the
editor, visualizes the resulting MIDI on a piano roll, and auditions the output with in-browser synthesis.
Version~2.0 extends this playground with account-based persistence: PostgreSQL stores users, sessions, and saved
source files; authenticated users get a file explorer, manual save with dirty-state tracking, and per-user document
tabs alongside bundled example compositions.

\subsection{Runtime topology}\label{subsec:studio-runtime-topology}

Production deployment is documented below.

\topic{Local development.} Uses Docker Compose with four application services plus PostgreSQL: database,
Express server (runs migrations then the API), Next.js client, and an nginx HTTPS gateway at
\texttt{motivo-studio.local}.
The gateway terminates TLS with a self-signed certificate, forwards page traffic to the client, and proxies
\texttt{/api/*} to the server with the \texttt{/api} prefix stripped, matching the URL shape the production client
exposes through Next.js rewrites.

\begin{figure}[htbp]
  \centering
  \motivoscalefig{%
    \begin{tikzpicture}[
        service/.style={motivo/service, minimum width=3.2cm, minimum height=1.05cm},
        store/.style={rounded corners=4pt, draw=motivoInk!50, thick, minimum width=3.0cm, minimum height=0.85cm,
        align=center, font=\scriptsize\sffamily},
        arrow/.style={motivo/arrow}
      ]
      \node[service, fill=motivoBlue!14] (browser) {Browser\\HTTPS};
      \node[service, fill=motivoBlue!12, below=0.75cm of browser] (client) {Next.js client\\Monaco, piano
      roll, auth UI};
      \node[service, fill=motivoCyan!16, below left=0.85cm and 0.15cm of client] (server) {Express
      API\\compile, auth, files};
      \node[service, fill=motivoGreen!16, below right=0.85cm and 0.15cm of server] (compiler)
      {\texttt{motivoc}\\subprocess};
      \node[store, fill=motivoSoft, below=0.7cm of server] (db) {PostgreSQL\\users, sessions, files};

      \draw[arrow] (browser) -- node[right, font=\scriptsize\sffamily, align=left] {pages +\\\texttt{/api/*}} (client);
      \draw[arrow] (client) -- node[above, font=\scriptsize\sffamily] {\texttt{/api/*} proxy} (server);
      \draw[arrow] (server) -- (compiler);
      \draw[arrow] (server) -- (db);
    \end{tikzpicture}%
  }
  \caption{Logical Motivo Studio topology shared by local Docker Compose and production Railway.
    The C++ compiler is not linked into the browser; the server spawns \texttt{motivoc} as a subprocess and reads
  temporary \texttt{.mid} output.}\label{fig:studio-architecture}
\end{figure}

The compile path deliberately keeps a narrow boundary: TypeScript in the client and server, C++ only inside
\texttt{motivoc}.
The server validates compile bodies with Zod, writes source to a temporary \texttt{.motivo} file, invokes
\texttt{execFile} on \texttt{COMPILER\_BIN} (default \texttt{/usr/local/bin/motivoc} in the server image), and returns
either \texttt{audio/midi} bytes or JSON diagnostics parsed from stderr.
File and auth routes live on the same Express app but require a valid session cookie.

\subsection{Backend API surface}\label{subsec:studio-api-surface}

Routes are mounted at the server root; the client always calls them under the \texttt{/api} prefix.

\begin{table}[htbp]
  \centering
  \renewcommand{\arraystretch}{1.25}
  \begin{tabularx}{\textwidth}{@{}M{3.4cm}>{\raggedright\arraybackslash\hspace{0pt}}X@{}}
    \toprule
    \textbf{Route group} & \textbf{Endpoints and behavior} \\
    \midrule
    Health & \texttt{GET /health}: liveness for orchestration. \\
    Compile & \texttt{POST /compile}: body \texttt{\{ source \}}; success returns MIDI; failure returns structured
    diagnostics (severity, \texttt{type}, message, optional line/column). \\
    Auth & \texttt{POST /auth/register}, \texttt{POST /auth/login}, \texttt{POST /auth/logout}, \texttt{GET /auth/me}:
    email/password credentials; HTTP-only session cookie (name from \texttt{SESSION\_COOKIE\_NAME}, secure flag
    enabled in production). \\
    Files & \texttt{GET /files}, \texttt{POST /files}, \texttt{GET /files/:id}, \texttt{PATCH /files/:id},
    \texttt{DELETE /files/:id}, \texttt{GET /files/:id/download}: all require authentication; store Motivo source
    per user in PostgreSQL. \\
    \bottomrule
  \end{tabularx}
  \caption{Express API groups implemented in the Motivo Studio server.}\label{tab:studio-api}
\end{table}

\subsection{Database schema}\label{subsec:studio-database}

PostgreSQL stores three entity groups defined in the server migration
\texttt{001\_auth\_account\_file\_storage\_schema}:

\begin{itemize}
  \item \textbf{users}: \texttt{id}, normalized \texttt{email}, \texttt{password\_hash}, optional
    \texttt{last\_opened\_file\_id}, timestamps.
  \item \textbf{motivo\_files}: per-user documents with \texttt{owner\_id} foreign key, \texttt{name},
    \texttt{source} text, \texttt{last\_opened\_at}, and unique \texttt{(owner\_id, lower(name))} constraint.
  \item \textbf{sessions}: \texttt{user\_id}, hashed \texttt{token\_hash}, \texttt{expires\_at}, indexed for lookup
    and expiry sweeps.
\end{itemize}

\subsection{Client feature modules}\label{subsec:studio-client-modules}

The Next.js client is organized by feature directories rather than a flat component tree.

\topic{Editor (\texttt{features/editor}).} Monaco with a registered Motivo language id, Monarch tokenizer
(keywords, types, instruments, drum notes, literals), light/dark themes, snippet completions, and marker mapping from
compiler diagnostics.

\topic{Compile (\texttt{features/compile}).} Axios client posting to \texttt{/api/compile} with
\texttt{arraybuffer} responses; distinguishes MIDI success from JSON diagnostic payloads.

\topic{MIDI and piano roll (\texttt{features/midi}, \texttt{features/piano-roll}).} Parse binary MIDI into
note events; canvas piano-roll with pitch scale, time axis, and playhead synchronization.

\topic{Playback (\texttt{features/playback}).} Tone.js scheduling with SoundFont-based General MIDI
instruments; transport bar for play, stop, progress, and MIDI download.

\topic{IDE shell (\texttt{features/ide}).} Workspace layout (editor pane, visualizer, logs, header);
\texttt{useIdeCompile} ties compile actions to editor markers and MIDI context; keyboard shortcuts for compile, save,
playback, and log panel focus.

\topic{Auth (\texttt{features/auth}).} React context for register, login, logout, and session refresh via
\texttt{/api/auth/*}.

\topic{Files (\texttt{features/files}).} File explorer, multi-tab documents (scratch buffer, bundled
examples, user files), manual save with dirty-state tracking, download of \texttt{.motivo} sources, create/rename/delete
for owned files.

\topic{Examples (\texttt{features/examples}).} Shipped \texttt{.motivo} sources
(e.g.\ \texttt{example.motivo},
\texttt{fur\_elise.motivo}) readable without login; opening an example does not require persistence.

Guests can compile and audition from the scratch editor and read-only examples; persistence and manual save require
authentication.
